% XAI Evaluation Card: blank LaTeX template
%
% The field descriptions are Table 1 of:
%   Rokas Gipiškis and Olga Kurasova. 2026. Evaluation Cards for XAI Metrics.
%   Proceedings of the Workshop on Evaluating Evaluations (EvalEval), 245--251.
%   https://aclanthology.org/2026.evaleval-1.39/
%
% Requires in the preamble: \usepackage{booktabs, tabularx, xcolor, url}
%
% Replace each \cardprompt{...} with your entry. For a single-column layout, use
% table instead of table* and \linewidth instead of \textwidth. A filled example
% is in examples/deletion-auc.tex.

% Remove this macro once every prompt has been replaced.
\providecommand{\cardprompt}[1]{\textit{\textcolor{gray}{#1}}}

\begin{table*}[t]
\centering
\small
\begin{tabularx}{\textwidth}{@{}l X@{}}
\toprule
\multicolumn{2}{c}{\textbf{XAI Evaluation Card}} \\
\midrule

\multicolumn{2}{@{}l}{\textbf{I. Identity}} \\[2pt]
Metric Name &
  \cardprompt{Unique, descriptive name for the evaluation metric.} \\[3pt]
Target Property / Properties &
  \cardprompt{List all explainability properties this metric operationalizes
  (e.g., fidelity, robustness, clarity), with references to definitions used.} \\[3pt]
Grounding Level &
  \cardprompt{One or more of: functionally-grounded / human-grounded /
  application-grounded (Doshi-Velez and Kim, 2017).} \\[3pt]
\addlinespace

\multicolumn{2}{@{}l}{\textbf{II. Scope and Context}} \\[2pt]
Evaluation Context &
  \cardprompt{Model architecture, data modality, and explanation scope
  (local / global) under which results are reported.} \\[3pt]
Assumptions &
  \cardprompt{All assumptions required by the metric (e.g., feature
  independence, locality, linearity, calibrated probabilities, meaningful
  baselines).} \\[3pt]
\addlinespace

\multicolumn{2}{@{}l}{\textbf{III. Implementation and Validation}} \\[2pt]
Implementation Available? &
  \cardprompt{Yes / No. If yes, provide URL or repository reference.} \\[3pt]
Validation Evidence &
  \cardprompt{Summary of sensitivity analysis, stability analysis, and
  correlation with related metrics. Report computational cost where relevant.} \\[3pt]
Gaming Risk &
  \cardprompt{How a method could achieve a high score on this metric without
  improving the target explainability property.} \\[3pt]
Known Failure Cases &
  \cardprompt{Conditions under which the metric is known to fail or produce
  misleading results.} \\[3pt]
\addlinespace

\multicolumn{2}{@{}l}{\textbf{IV. Relationships and Limitations}} \\[2pt]
Relationship to Other Metrics &
  \cardprompt{Metrics targeting the same property. Known agreements or
  disagreements in results.} \\[3pt]
Disagreement Handling &
  \cardprompt{If this metric conflicts with others reported, state which
  property is prioritised for the target deployment scenario and why.} \\[3pt]
Limitations &
  \cardprompt{Main limitations as an operationalization of the target property.
  Note contexts where the metric should not be used.} \\[3pt]

\bottomrule
\end{tabularx}
\caption{XAI Evaluation Card for \cardprompt{metric name}. Fields marked N/A
require a brief justification. Template from~\cite{gipiskis-kurasova-2026-evaluation}.}
\label{tab:xai-evaluation-card}
\end{table*}

% BibTeX:
%
% @inproceedings{gipiskis-kurasova-2026-evaluation,
%     title = "Evaluation Cards for {XAI} Metrics",
%     author = "Gipi{\v{s}}kis, Rokas  and
%       Kurasova, Olga",
%     editor = "Akhtar, Mubashara  and
%       Batzner, Jan  and
%       Choshen, Leshem  and
%       Ghosh, Avijit  and
%       Gohar, Usman  and
%       Mickel, Jennifer  and
%       Pant, Ichhya  and
%       Talat, Zeerak  and
%       Lin, Michelle",
%     booktitle = "Proceedings of the Workshop on Evaluating Evaluations ({E}val{E}val)",
%     month = jul,
%     year = "2026",
%     address = "San Diego, CA",
%     publisher = "Association for Computational Linguistics",
%     url = "https://aclanthology.org/2026.evaleval-1.39/",
%     doi = "10.18653/v1/2026.evaleval-1.39",
%     pages = "245--251",
%     ISBN = "979-8-89176-429-3"
% }
